% 独立章节文件：可单独编辑，统一由 main.tex 引入。
\chapter{自动温补：实验设计与运行流程}
\section{两类设备在温补阶段的分工}
常温 2TD600 与温控 2FDT600WK 是不同驱动路径。温补实验只允许选用\textbf{2FDT600WK}，同时复用经过现场验证的原 PLC 温控驱动；IMU 原始帧通过同一实时数据队列采集。\textbf{温箱设定值 SP 不是样品温度}，模型使用同步获取的 IMU 芯片温度 $T_c$。

\section{默认温度计划：拟合点与温度验证点隔离}
% 可复用表格：修改内容后重新编译 main.tex
\begin{table}[H]\centering
\caption{软件默认的五温度点自动实验计划}
\begin{tabular}{@{}cccp{55mm}@{}}\toprule
次序 & 目标 SP & 数据角色 & 实验目的\\\midrule
1 & $-20\, ^\circ\mathrm C$ & FIT & 低温参数拟合\\
2 & $0\, ^\circ\mathrm C$ & VAL & 独立温度验证（不得用于拟合）\\
3 & $25\, ^\circ\mathrm C$ & FIT & 常温温度模型锚定\\
4 & $45\, ^\circ\mathrm C$ & VAL & 独立中高温验证\\
5 & $70\, ^\circ\mathrm C$ & FIT & 高温参数拟合\\\bottomrule\end{tabular}\end{table}

程序要求至少五个互异温度点，升序输入，允许的设定范围为 $-40\sim85\, ^\circ\mathrm C$；数据角色采用\textbf{第 1、3、5… 点为拟合，第 2、4… 点为验证}的交替规则。增加温度点应重新审查此规则，以避免将本应验证的温度误归为拟合温度。一次运行只覆盖\textbf{单向升温}；成功后默认回到 $25\, ^\circ\mathrm C$，但回温不执行完整的参数复测，不能宣称已完成冷热滞回验证。

\section{完整任务状态机}
\begin{figure}[H]\centering
\begin{tikzpicture}[>=Stealth,node distance=4.6mm, every node/.style={font=\footnotesize},task/.style={draw=border,fill=soft,rounded corners=1.5mm,minimum height=8mm,text width=108mm,align=center}]
\node[task] (a) {设备与安全预检：2FDT、IMU、PLC、使能、限幅、坐标、配置锁定};
\node[task,below=of a] (b) {PLC 设置定值模式 $\rightarrow$ 写 SP $\rightarrow$ 启动 $\rightarrow$ SP/运行状态回读确认};
\node[task,below=of b] (c) {等 PV 在目标 $\pm1\, ^\circ\mathrm C$ 内连续 30 s；无报警且反馈新鲜};
\node[task,below=of c] (d) {至少 30 min 保温 + 芯片温度 60 s 波动不大于 $0.2\, ^\circ\mathrm C$};
\node[task,below=of d] (e) {ACC 24 姿态 + GYRO 零速；按选项加入 GYRO 正反转动态工况};
\node[task,below=of e] (f) {逐工况质量检查、重试、停车反馈、记录原始帧和汇总};
\node[task,below=of f] (g) {下一个温度点；全部采完后回 $25\, ^\circ\mathrm C$ 并验证 PV};
\node[task,below=of g] (h) {温度拟合、独立温度验证、对照 RMSE、候选 JSON / Excel};
\foreach \x/\y in {a/b,b/c,c/d,d/e,e/f,f/g,g/h}{\draw[->,thick,blue] (\x.south)--(\y.north);}
\node[draw=amber,fill=amber!6,rounded corners=1mm,align=center,text width=36mm,right=5mm of d,font=\scriptsize] (stop) {任何报警/超时/数据断流\\受控中止 + 停车确认};
\draw[->,thick,amber,dashed] (d.east)--(stop.west);
\end{tikzpicture}
\caption{自动温补状态机；上图为工程逻辑示意，现场安全链路必须独立于 UI。}
\end{figure}


\section{PLC 自动控温的当前判据}
% 可复用表格：修改内容后重新编译 main.tex
\begin{table}[H]\centering\small
\caption{PLC 自动流程当前代码判据}
\noindent\begin{tabularx}{\linewidth}{@{}Yp{41mm}Y@{}}\toprule
项目 & 当前值 & 目的\\\midrule
PV/SP/状态/报警等更新最大年龄 & 15 s & 禁止把缓存温度当成有效反馈\\
SP 写入后回读允许误差 & $\pm0.2\, ^\circ\mathrm C$ & 检查设温命令确已应用\\
模式与运行 & 定值模式 + 已启动 + 运行中 & 仅 SP 一致不足以开始试验\\
SP 回读/启动确认超时 & 90 s & 拒绝控制器无响应\\
到温判定 & PV 距目标 $\le1\, ^\circ\mathrm C$ 连续 30 s & 初步到温判定\\
温箱到温超时 & 6 h & 避免无限等待\\
PLC 报警/故障/断连 & 立即中止 & 防止失控状态下旋转和采样\\\bottomrule\end{tabularx}\end{table}

数据通过原\Code{tc_set_fixed_value_mode()}、\Code{tc_sp_set()}、\Code{tc_fixed_value_startup()} 等接口操作。已通过现场使用事实确认原上位机曾成功读 PV 与设 SP；但是\textbf{新增自动循环程序本身尚未完成该设备的连续实机联动验收}。

\section{芯片热平衡判据（软件初值）}
% 可复用表格：修改内容后重新编译 main.tex
\begin{table}[H]\centering\small
\caption{IMU 内部温度稳定性要求}
\noindent\begin{tabularx}{\linewidth}{@{}Yp{38mm}Y@{}}\toprule
条件 & 默认值 & 技术说明\\\midrule
到温后最短保温 & 1,800 s（30 min） & 不是只看温箱面板达到 SP\\
最近 60 s 芯片温度波动 & $\le0.2\, ^\circ\mathrm C$ & 当前至少积累 20 个稳定检查样本\\
芯片温度与 PV 差 & $\le10\, ^\circ\mathrm C$ & 宽松保护门限，\textbf{不等于计量级热平衡}\\
热平衡等待超时 & 6 h & 超时中止\\
单工况芯片温度跨度 & $\le0.5\, ^\circ\mathrm C$ & 温补任务的临时加强判据\\\bottomrule\end{tabularx}\end{table}

建议正式研发记录另加 PCB 表面参考温度、箱内温度空间差和保温时长试验，以确定 30 min 是否足够；不得单凭芯片与箱温相差 $\le10\, ^\circ\mathrm C$ 宣称绝对热平衡。

\section{各温度点的 ACC/GYRO 工况数量}
% 可复用表格：修改内容后重新编译 main.tex
\begin{table}[H]\centering
\caption{一个温度点的任务量}
\noindent\begin{tabularx}{\linewidth}{@{}YccY@{}}\toprule
工作项 & FIT & VAL & 说明\\\midrule
ACC 静态多姿态 & 12 & 12 & 按常温姿态表执行\\
GYRO 零速 & 1 & 1 & 拟合零速按配置 120 s，验证至少 10 s\\
GYRO 动态（可选） & 18 & 12 & $\pm1,\pm2,\pm5$ 拟合；$\pm1.5,\pm3.5$ 验证\\\midrule
关闭动态 GYRO & \multicolumn{2}{c}{\textbf{26 工况}} & 仅可生成 GYRO 零偏温补，不能声称完成比例因子温补\\
启用动态 GYRO & \multicolumn{2}{c}{\textbf{56 工况}} & 5 个温度点总计 280 工况\\\bottomrule\end{tabularx}\end{table}

每个工况都需要转台实际角度/速率、芯片温度和原始 IMU 帧。动态速率必须先通过反馈稳定检查；当前 2FDT 工程的速率反馈依据编码器位置差分估计，\textbf{尚须通过参考计量证明速度不确定度满足目标}。
